\chapter{The claim inventory and the banned claims}
\label{ch:claims}

\section*{What this chapter is for}

By the time you apply you will have written the same facts in four places: a
curriculum vitae, a covering letter, a set of form answers, and the README of
whatever you built. They will not agree, because they were written on different
days for different purposes, and the disagreements are the cheapest thing in
your application for somebody to find.

This chapter is two routines. One makes the four documents agree. The other
stops you saying, out loud and under pressure, something your own work does not
support.

\section{The claim inventory}

Every sentence in an application that could be checked is a claim. Collect them
into one table.

\begin{kittable}{@{}lXl@{}}
\textbf{Column} & \textbf{Holds} & \textbf{Why} \\
\midrule
Claim & the sentence, as written & the thing that will be read \\
Where & every document it appears in & one edit has to reach all of them \\
Backed by & the file, commit, measurement or person & if this is blank it is not a claim, it is a hope \\
Strength & measured / observed / reported / judgement & tells you which sentence to soften \\
Last checked & a date & an unchecked claim is a stale claim waiting \\
\end{kittable}

\code{templates/claim-inventory.md} is this table. Filling it the first time
takes about an hour and reliably finds three things: a number that appears at
two different precisions in two documents, a claim whose backing turns out to
be a memory, and a capability described more strongly in the covering letter
than in the repository it refers to.

\begin{kitnote}
The \emph{strength} column is the one that changes what you write. A claim
backed by a measurement can be stated flatly. One backed by an observation on
one occasion needs the occasion in the sentence. One backed by judgement needs
to be marked as judgement \dash{} which is exactly the distinction this book
makes on its own pages, for the same reason.
\end{kitnote}

\section{Why inconsistency costs more than it should}

A reader who finds one discrepancy does not discount one claim. They lose their
basis for accepting any of the others without checking, and they cannot check
most of them.

\judgement{This is why consistency is worth an hour even though no single
inconsistency is serious. You are not defending each sentence; you are
defending the reader's willingness to take the unverifiable ones on trust, and
that is a single asset that fails all at once.}

The corollary is practical: a claim you cannot back should be \emph{removed},
not softened. A softened unbacked claim is still unbacked and now reads as
evasive.

\section{The banned-claims list}

The inverse document, and the one that matters most in the room.

Before any conversation, write down the sentences you must not say: the things
that would be convenient, that you half-believe, and that your evidence does
not actually support. Keep it to one page. Read it before the call.

Typical entries, and each is a real shape rather than a hypothetical:

\begin{itemize}
  \item A capability your integration has never exercised against the live
        thing \dash{} you built it against a specification and it has not run.
  \item A count you cannot reproduce on the day: how many tests, how many
        scenarios, how much something improved.
  \item An outcome you contributed to and did not own, stated as though you
        owned it.
  \item A superlative \dash{} only, first, every time, always. These are the
        most dangerous entries on the list, because the interviewer needs one
        counterexample and often has one.
  \item Work by somebody else that you were near.
\end{itemize}

\begin{kittrap}
The single most expensive sentence available to you is a superlative about your
own checks. ``It caught every one'' is a sentence people say about their own
test suites because it is what they remember, and it is very often false by one
case. If one of your deliberately broken variants survived, that is a better
story than the one you were about to tell: it is the story of an instrument
being tested and found imperfect, which is what rigour looks like from the
outside.
\end{kittrap}

\section{The recovery line}

Somebody will eventually find something in your repository that has stopped
being true, or ask for a number you cannot produce from memory. Prepare the
sentence now, because the unprepared version is worse than the problem.

\begin{kitmethod}
\textbf{For a stale claim:} name it as stale, say when it was true, and say what
it is now if you know \dash{} or say that you will check rather than guess. Do
not defend it and do not apologise at length. One sentence, then continue.

\textbf{For a number you cannot reproduce:} say that you will not give it from
memory, name the file it is in, and give the shape instead \dash{} the order of
magnitude, or the direction. \judgement{Refusing to produce a number you are
not sure of is read as discipline, every time, by anybody who has been burnt by
somebody who did not refuse.}
\end{kitmethod}

\begin{drill}
Take the artefact you would send today and read its README against the
repository, line by line, with the inventory table open. Note every sentence
that has stopped being exactly true.

Then write the banned-claims page. The first entry is usually the most
impressive-sounding sentence you have been using, and finding that out on a
Tuesday is considerably cheaper than finding it out in the room.
\end{drill}

\section*{What this chapter does not claim}

It does not claim interviewers routinely audit portfolios against applications.
Most do not have time. The claim is that the cost of being caught once is
disproportionate, and the cost of prevention is an hour.

The recovery lines are the author's judgement about what reads well under
pressure. They are stated as judgement and they are not measured.
